\subsection{Boty heurystyczne}

Boty heurystyczne rozwijane były iteracyjnie: każda kolejna wersja rozszerza poprzednią o nowe elementy strategii. Dzięki temu można ocenić wpływ poszczególnych mechanizmów na skuteczność oraz obserwować, jak nawet proste heurystyki poprawiają decyzje względem losowego wyboru akcji.

\subsubsection{it1 -- priorytet punktów zwycięstwa}

It1 wybiera akcje według \textbf{stałej hierarchii priorytetów}: najpierw budowa miasta, potem osady, drogi, karty rozwoju, na końcu handel i inne. W obrębie każdej kategorii wybór jest \textbf{losowy}. Bot nie handluje z bankiem, nie optymalizuje ustawień początkowych ani rozbójnika. Strategia jest prosta i krótkowzroczna --- maksymalizuje natychmiastowy przyrost punktów, nie uwzględniając przyszłej produkcji ani pozycji na planszy.

\subsubsection{it2 -- handel z bankiem i cel zakupu}

It2 rozszerza it1 o \textbf{aktywne wykorzystanie handlu z bankiem}. Wykrywa kontrolowane porty (2:1, 3:1, 4:1) i ocenia, \textbf{do którego celu zakupu jest najbliżej po uwzględnieniu możliwych wymian} --- tzn.\ ile brakujących surowców pozostaje po optymalnym ``wydaniu'' nadwyżek. Cel wybierany jest nie na podstawie tego, co da się kupić od razu, lecz tego, do czego po handlu jest się najbliżej. Bot może odłożyć budowę drogi i najpierw handlować, jeśli po wymianie będzie mógł zbudować osadę lub miasto. Hierarchia: miasto i osada (natychmiast) $\rightarrow$ handel w kierunku wybranego celu $\rightarrow$ droga $\rightarrow$ karta rozwoju $\rightarrow$ koniec tury.

\subsubsection{it3 -- ustawienia początkowe i rozbójnik}

It3 dodaje \textbf{świadomą fazę początkową} oraz \textbf{celowe ustawianie rozbójnika}. Ustawienia początkowe oceniane są według produkcji (z naciskiem na surowce ważne we wczesnej grze), różnorodności zasobów i dostępu do portów. Drugie ustawienie dobierane jest tak, aby \textbf{uzupełniać braki} po pierwszym --- bot dąży do pokrycia wszystkich typów surowców. Przy ruchu rozbójnikiem bot wybiera heks, który \textbf{maksymalnie ogranicza produkcję przeciwnika} przy \textbf{minimalnym wpływie na własną}. W fazie głównej zachowuje logikę it2 (handel i cele zakupu).

\subsubsection{it4 -- karty rozwoju i deterministyczny wybór budowy}

It4 wprowadza \textbf{inteligentne używanie kart rozwoju} oraz \textbf{deterministyczny wybór lokalizacji} zamiast losowego rozstrzygania remisów. Dla każdego typu karty (Rycerz, Monopol, Rozwój itd.) stosowana jest odrębna logika: Rycerze --- blokada przeciwnika i premia Największa Armia, karty zasobowe --- pod kątem odblokowania budowy miasta lub osady, karty VP --- jako bezpośredni zysk punktowy. Uwzględniane jest \textbf{ryzyko przed rzutem} (7): unika się akcji, które mocno powiększają rękę przed rzutem. Budowa (miasto, osada, droga) wybierana jest deterministycznie na podstawie \textbf{potencjału produkcyjnego} i użyteczności sieci dróg; priorytety dostosowują się do fazy gry.

\subsubsection{it5 -- ocena pozycji i jednokrokowy lookahead}

It5 uzupełnia it4 o \textbf{jednokrokową symulację} (\emph{one-step lookahead}) i \textbf{zunifikowaną ocenę pozycji}. Zamiast sztywnych priorytetów bot dla każdej rozważanej akcji deterministycznej: stosuje ją do kopii stanu, liczy wartość funkcji oceny pozycji, cofa akcję --- i wybiera akcję dającą \textbf{najwyższą ocenę}. Decyzja opiera się więc na \textbf{rzeczywistym wpływie na stan gry}, a nie na samym typie akcji. Najpierw rozważane są budowy miasta i osady (wybór najlepszej lokalizacji), potem drogi i handel; dodatkowo premiowane są ruchy, które \textbf{odblokowują} w następnym kroku budowę miasta lub osady.

Zakup karty rozwoju nie jest symulowany (losowość talii) i oceniany heurystycznie; gdy najlepszą opcją jest zakończenie tury bez poprawy, bot \textbf{zawraca do logiki it4}. Przy wyrzuceniu 7 it5 \textbf{nie odrzuca losowo}: ustala najbliższy cel (miasto, osada, droga, karta), po czym odrzuca surowce najmniej potrzebne do tego celu, chroniąc zasoby krytyczne. It5 łączy deterministyczny wybór i sensowne priorytety z oceną konsekwencji ruchu przy niewielkim koszcie obliczeniowym.
